\documentclass[11pt]{article}
\newcommand{\DocTitle}{Twelve-Week Engineering Curriculum}
\newcommand{\DocSubtitle}{Topologically ordered skills, labs, capstone and assessments}
\newcommand{\DocID}{DOS-06}
% design-of-sites-style.tex
% Suite preamble for the Web Design Patterns / Design of Sites suite.
%
% RECONSTRUCTION PROVENANCE: the nine uploaded suite documents load a shared
% preamble with `\input{style.tex}`, but no `style.tex` was present in the
% upload batches or anywhere in this repository's history. This file is a
% documented, suite-specific reconstruction inferred from the interface the
% documents actually use: per-document metadata (\DocTitle, \DocSubtitle,
% \DocID defined before this file is input), the front-matter renderer
% \MakeFront, the provenance marker \Provenance, the callout environments
% groundedbox/modernbox/notebox/decisionbox, and the colors
% Ink/Ocean/Mint/Amber/Faint referenced in document TikZ code. It is NOT the
% recovered original style; the original visual design is unknown.
%
% pdfLaTeX only.
\usepackage[T1]{fontenc}
\usepackage[utf8]{inputenc}
\usepackage{lmodern}
\usepackage{microtype}
\usepackage[letterpaper,margin=0.85in,headheight=15pt]{geometry}
\usepackage[table]{xcolor}
\definecolor{Ink}{HTML}{1E2A36}
\definecolor{Ocean}{HTML}{0F4C72}
\definecolor{Mint}{HTML}{0E7E6B}
\definecolor{Amber}{HTML}{B07A1E}
\definecolor{Faint}{HTML}{F2F7F9}
\usepackage{amsmath,amssymb}
\usepackage{array,booktabs,longtable,tabularx,multirow}
\usepackage{enumitem}
\usepackage[most]{tcolorbox}
\usepackage{listings}
\usepackage{tikz}
\usetikzlibrary{arrows.meta,positioning,shapes.geometric,fit,calc}
\usepackage{needspace}
\usepackage{titlesec}
\usepackage{fancyhdr}
\usepackage{lastpage}
\usepackage{xurl}
\usepackage{hyperref}
\hypersetup{colorlinks=true,linkcolor=Ocean,urlcolor=Mint,citecolor=Ocean,pdftitle={\DocTitle},pdfauthor={Web Design Patterns Engineering Suite}}
\setlength{\parindent}{0pt}
\setlength{\parskip}{5pt plus 1pt}
\setlength{\emergencystretch}{2em}
\renewcommand{\arraystretch}{1.18}
\setlist[itemize]{itemsep=2pt,topsep=3pt,leftmargin=1.5em}
\setlist[enumerate]{itemsep=2pt,topsep=3pt,leftmargin=1.7em}
\titleformat{\section}{\Large\bfseries\color{Ocean}}{\thesection}{0.7em}{}[\vspace{-3pt}\color{Mint}\titlerule]
\titleformat{\subsection}{\large\bfseries\color{Ink}}{\thesubsection}{0.6em}{}
\titleformat{\subsubsection}{\normalsize\bfseries\color{Mint}}{\thesubsubsection}{0.55em}{}
\pagestyle{fancy}
\fancyhf{}
\fancyhead[L]{\sffamily\footnotesize\color{Ocean} Design of Sites | \DocID}
\fancyhead[R]{\sffamily\footnotesize\color{Ink} Web Design Patterns Suite}
\fancyfoot[L]{\sffamily\scriptsize Original interpretive guidance; pattern names used as a taxonomy}
\fancyfoot[R]{\sffamily\scriptsize\thepage\ / \pageref{LastPage}}
\renewcommand{\headrulewidth}{0.3pt}
\lstdefinestyle{dositescode}{basicstyle=\ttfamily\footnotesize,breaklines=true,columns=fullflexible,keepspaces=true,showstringspaces=false,frame=single,rulecolor=\color{Mint},backgroundcolor=\color{Faint},aboveskip=8pt,belowskip=8pt}
\lstset{style=dositescode}
% Callouts. groundedbox marks statements verified against the public
% bibliographic record; modernbox marks newly authored modern engineering
% interpretation; notebox is a study/summary callout; decisionbox records
% decisions and trade-offs.
\newtcolorbox{groundedbox}[1][]{enhanced,breakable,colback=Faint,colframe=Mint,boxrule=0.5pt,arc=1.5mm,title={Grounded in the public record},fonttitle=\bfseries\small,#1}
\newtcolorbox{modernbox}[1][]{enhanced,breakable,colback=white,colframe=Ocean,boxrule=0.5pt,arc=1.5mm,title={Modern engineering interpretation},fonttitle=\bfseries\small,#1}
\newtcolorbox{notebox}[1][]{enhanced,breakable,colback=Faint,colframe=Amber,boxrule=0.5pt,arc=1.5mm,#1}
\newtcolorbox{decisionbox}[1][]{enhanced,breakable,colback=white,colframe=Amber,boxrule=0.5pt,arc=1.5mm,title={Decision record},fonttitle=\bfseries\small,#1}
% Front matter renderer: consumes \DocTitle, \DocSubtitle, \DocID.
\newcommand{\MakeFront}{%
  \thispagestyle{empty}%
  {\color{Mint}\rule{\textwidth}{3pt}}\par\vspace{16pt}
  {\sffamily\small\color{Amber}\bfseries \DocID\par}\vspace{10pt}
  {\Huge\bfseries\color{Ocean} \DocTitle\par}\vspace{12pt}
  {\Large\color{Ink} \DocSubtitle\par}\vspace{14pt}
  {\small\color{Ink} Independent educational and engineering companion to \emph{The Design of Sites} (2nd edition). Original explanatory material; the copyrighted pattern text is not reproduced.\par}
  \vspace{10pt}{\color{Mint}\rule{\textwidth}{1pt}}\par\vspace{8pt}
  \tableofcontents\clearpage}
% Provenance marker used at the top of each document body.
\newcommand{\Provenance}{\par\noindent{\footnotesize\sffamily\color{Amber}\bfseries Provenance:}~{\footnotesize\color{Ink}\itshape claims grounded in the public bibliographic record are boxed and cited; everything else is newly authored interpretation.}\par\smallskip}

\begin{document}
\MakeFront
\section{Learning architecture}

\Provenance
\begin{groundedbox}The source mapping uses the historical chapter and pattern group identifiers [B1]. The 12-week order, prerequisites, lab assignments and grading criteria are original and are not the book's teaching sequence.\end{groundedbox}
\textbf{Format.} Twelve weekly modules, approximately four to seven guided hours plus individual practice. The same reference application evolves from a navigable prototype into an operational service. Learners should work in pairs for accessibility observation and peer review.
\textbf{Baseline skills.} HTML, CSS, JavaScript or TypeScript, HTTP, elementary APIs, version control and test automation. Additional support may be needed for threat modeling and cloud deployment.

\section{Prerequisite map}

\begin{notebox}\textbf{Critical path:} user research $\rightarrow$ information architecture $\rightarrow$ semantic interaction $\rightarrow$ responsive patterns $\rightarrow$ security and domain workflows $\rightarrow$ architecture $\rightarrow$ observability $\rightarrow$ delivery assurance $\rightarrow$ capstone.

\textbf{Parallel path:} content/search after information architecture; recommendation moderation after search and identity. Avoid complex service decomposition until a complete user journey works in the modular monolith.
\end{notebox}

\section{Week 01: Evidence and user missions}
\textbf{Prerequisites:} None.\\
\textbf{Pattern focus:} A1--A12; foundation 1/3.\\
\textbf{Hands-on lab:} Select a scenario, interview stakeholders and create task inventory..\\
\textbf{Required artifacts:} Personas-as-hypotheses, top tasks, research consent, user outcome baseline.\\
\textbf{Acceptance criteria:} Task goals grounded; assumptions tagged; privacy risks identified.
\begin{decisionbox}\textbf{Weekly reflection:} Which customer outcome improved? Which test could falsify the claimed improvement? What was introduced as architectural or privacy risk?\end{decisionbox}
\section{Week 02: Information architecture}
\textbf{Prerequisites:} Week 1.\\
\textbf{Pattern focus:} B1--B9; C1--C2.\\
\textbf{Hands-on lab:} Design taxonomy, sitemap and dual task/category entry paths..\\
\textbf{Required artifacts:} Tree test, route map, homepage hypothesis.\\
\textbf{Acceptance criteria:} At least five task paths tested; duplicate destinations resolved.
\begin{decisionbox}\textbf{Weekly reflection:} Which customer outcome improved? Which test could falsify the claimed improvement? What was introduced as architectural or privacy risk?\end{decisionbox}
\section{Week 03: Content and search}
\textbf{Prerequisites:} Week 2.\\
\textbf{Pattern focus:} D1--D11; J1--J3.\\
\textbf{Hands-on lab:} Define content schemas and build an ACL-safe search prototype..\\
\textbf{Required artifacts:} CMS type schema, content style guide, search evaluation corpus.\\
\textbf{Acceptance criteria:} Unauthorized items excluded; zero-results handled.
\begin{decisionbox}\textbf{Weekly reflection:} Which customer outcome improved? Which test could falsify the claimed improvement? What was introduced as architectural or privacy risk?\end{decisionbox}
\section{Week 04: Semantic interaction}
\textbf{Prerequisites:} Weeks 2--3.\\
\textbf{Pattern focus:} H6--H13; K1--K17.\\
\textbf{Hands-on lab:} Implement accessible navigation, form validation and one modal/dialog..\\
\textbf{Required artifacts:} Component stories, keyboard tests, screen-reader notes.\\
\textbf{Acceptance criteria:} No keyboard trap, clear error recovery, landmark map.
\begin{decisionbox}\textbf{Weekly reflection:} Which customer outcome improved? Which test could falsify the claimed improvement? What was introduced as architectural or privacy risk?\end{decisionbox}
\section{Week 05: Responsive design}
\textbf{Prerequisites:} Week 4.\\
\textbf{Pattern focus:} I1--I6; M1--M3.\\
\textbf{Hands-on lab:} Build mobile-to-desktop layout, print and location fallback..\\
\textbf{Required artifacts:} Responsive QA matrix and device evidence.\\
\textbf{Acceptance criteria:} Reflow at 320px, zoom and denied-permission path.
\begin{decisionbox}\textbf{Weekly reflection:} Which customer outcome improved? Which test could falsify the claimed improvement? What was introduced as architectural or privacy risk?\end{decisionbox}
\section{Week 06: Trust, identity and privacy}
\textbf{Prerequisites:} Weeks 1--5.\\
\textbf{Pattern focus:} E1--E9; H2--H5.\\
\textbf{Hands-on lab:} Integrate a demo identity provider and configurable privacy preferences..\\
\textbf{Required artifacts:} Threat model, session test plan, privacy flow diagram.\\
\textbf{Acceptance criteria:} Authz negative tests and consent withdrawal pass.
\begin{decisionbox}\textbf{Weekly reflection:} Which customer outcome improved? Which test could falsify the claimed improvement? What was introduced as architectural or privacy risk?\end{decisionbox}
\section{Week 07: Commerce state machines}
\textbf{Prerequisites:} Week 6.\\
\textbf{Pattern focus:} F1--F9; H1.\\
\textbf{Hands-on lab:} Model carts, checkout, payment callbacks and stock changes..\\
\textbf{Required artifacts:} State diagram, OpenAPI contract and transaction tests.\\
\textbf{Acceptance criteria:} Idempotent duplicate requests; no client-price trust.
\begin{decisionbox}\textbf{Weekly reflection:} Which customer outcome improved? Which test could falsify the claimed improvement? What was introduced as architectural or privacy risk?\end{decisionbox}
\section{Week 08: Relevance and community}
\textbf{Prerequisites:} Weeks 3, 6.\\
\textbf{Pattern focus:} G1--G7; A3.\\
\textbf{Hands-on lab:} Add opt-in recommendations or moderated reviews in a sandbox..\\
\textbf{Required artifacts:} Bias and abuse cases, moderation policy, fallback.\\
\textbf{Acceptance criteria:} Opt-out honored and fake-review abuse test executed.
\begin{decisionbox}\textbf{Weekly reflection:} Which customer outcome improved? Which test could falsify the claimed improvement? What was introduced as architectural or privacy risk?\end{decisionbox}
\section{Week 09: Domain decomposition}
\textbf{Prerequisites:} Weeks 3, 7.\\
\textbf{Pattern focus:} A/B/F/H.\\
\textbf{Hands-on lab:} Identify bounded contexts, API owners and shared UI contracts..\\
\textbf{Required artifacts:} C4-style context/container diagrams and ADRs.\\
\textbf{Acceptance criteria:} Consistency boundary defended; contracts versioned.
\begin{decisionbox}\textbf{Weekly reflection:} Which customer outcome improved? Which test could falsify the claimed improvement? What was introduced as architectural or privacy risk?\end{decisionbox}
\section{Week 10: Performance and observability}
\textbf{Prerequisites:} Weeks 5, 9.\\
\textbf{Pattern focus:} L1--L6.\\
\textbf{Hands-on lab:} Instrument RUM, traces, backend latency and image budgets..\\
\textbf{Required artifacts:} Baseline/after graphs and SLO proposal.\\
\textbf{Acceptance criteria:} Published Web Vitals targets assessed with p75 data.
\begin{decisionbox}\textbf{Weekly reflection:} Which customer outcome improved? Which test could falsify the claimed improvement? What was introduced as architectural or privacy risk?\end{decisionbox}
\section{Week 11: Delivery and hardening}
\textbf{Prerequisites:} Weeks 4, 6, 10.\\
\textbf{Pattern focus:} B/E/H/J/K/L.\\
\textbf{Hands-on lab:} Introduce secret scans, SAST, API tests, A11Y and performance checks..\\
\textbf{Required artifacts:} CI gate matrix and remediation report.\\
\textbf{Acceptance criteria:} Critical journey, authz and keyboard gates pass.
\begin{decisionbox}\textbf{Weekly reflection:} Which customer outcome improved? Which test could falsify the claimed improvement? What was introduced as architectural or privacy risk?\end{decisionbox}
\section{Week 12: Capstone and review}
\textbf{Prerequisites:} Weeks 1--11.\\
\textbf{Pattern focus:} All groups.\\
\textbf{Hands-on lab:} Demo a resilient end-to-end flow with failure injection and research replay..\\
\textbf{Required artifacts:} Capstone portfolio, release notes, ADRs, usability evidence.\\
\textbf{Acceptance criteria:} Stakeholder review, rollback trial, residual risk sign-off.
\begin{decisionbox}\textbf{Weekly reflection:} Which customer outcome improved? Which test could falsify the claimed improvement? What was introduced as architectural or privacy risk?\end{decisionbox}
\section{Assessment rubric}

\begin{tabularx}{\linewidth}{@{}p{45mm}p{19mm}X@{}}\toprule
Competency & Weight & Exemplary evidence\\\midrule
User and domain reasoning & 20\% & Representative tasks, research evidence, explicit bounded contexts\\
Semantic UX and A11Y & 20\% & Keyboard/screen-reader tests, coherent navigation and error recovery\\
Security and privacy & 20\% & Threat cases, server-side authorization, controlled consent and logging\\
Architecture and contracts & 15\% & Clear ADRs, API/state contracts, independent concerns\\
Performance/operations & 15\% & Field metrics, tracing, failure testing and rollback\\
Technical communication & 10\% & Reproducible instructions, issue-to-pattern traceability\\
\bottomrule\end{tabularx}
\textbf{Pass criterion (program recommendation):} at least 80/100, with no unmitigated critical data exposure, loss-of-funds or fundamental keyboard inaccessibility in the final journey. Scores are program-designed, not from the source book.

\section{Capstone scenarios}
\begin{itemize}
\item 	extbf{Public service portal:} eligibility, identity, assisted completion, appointment status and privacy controls.
\item 	extbf{Small e-commerce storefront:} search, product, cart, idempotent checkout, fulfillment timeline and returns.
\item 	extbf{Knowledge community:} editorial publishing, search, moderated discussion, accessibility and personalization opt-out.
\item 	extbf{Internal service desk:} role-based navigation, searchable help, task forms and auditable workflow status.
\end{itemize}

\section{Instructor and peer-review protocol}

\begin{enumerate}
\item Before the module, verify that tools and linked specifications still have the versions recorded in [DOS-09].
\item Review artifacts against observable tests rather than framework preference.
\item Ask learners to distinguish the historic pattern title from their newly proposed implementation.
\item Run one adversarial test per module: malicious actor, assistive-technology user or partial outage.
\item Require a small ADR whenever the team departs from a simpler design, especially for runtime microfrontend composition.
\end{enumerate}


\clearpage
% design-of-sites-references.tex
% Shared references section for the Web Design Patterns / Design of Sites
% suite.
%
% RECONSTRUCTION PROVENANCE: the nine uploaded suite documents end with
% `\input{references.tex}`, but no `references.tex` was present in the upload
% batches or anywhere in this repository's history. The documents cite only
% keys [B1] and [B2] (bracketed tokens such as "[B4, F1, H10, ...]" are
% pattern-ID template placeholders, not citations). The two entries below are
% reconstructed from the bibliographic facts the documents themselves state:
% publisher record and bibliographic-database record of the second edition.
% This file is NOT the recovered original reference list.
\clearpage
\section*{References}
\begin{description}[leftmargin=3.2em,style=nextline]
  \item[{[B1]}] Publisher bibliographic record and public table of contents:
    D.~K.~van Duyne, J.~A.~Landay, J.~I.~Hong, \emph{The Design of Sites:
    Patterns, Principles, and Processes for Crafting a Customer-Centered Web
    Experience}, 2nd edition, Prentice Hall, December 2006, 1{,}024 pages.
    Confirms the five foundational chapters, the thirteen pattern groups
    (A--M), the 107 pattern labels and the research appendixes used as this
    suite's taxonomy.
  \item[{[B2]}] Bibliographic-database record (O'Reilly learning-platform
    catalog entry) for the same second edition; corroborates edition,
    authors, publication date and chapter/group structure. The copyrighted
    pattern explanations themselves were not retrieved or reproduced for
    this suite.
\end{description}

\end{document}
